
\documentclass[11pt]{article}

\usepackage[margin=1in]{geometry}
\usepackage{amsmath,amssymb,amsthm,mathtools}
\usepackage{enumitem}
\usepackage{hyperref}
\usepackage{mathrsfs}
\usepackage{microtype}

\hypersetup{
  colorlinks=true,
  linkcolor=blue,
  citecolor=blue,
  urlcolor=blue
}

\newtheorem{theorem}{Theorem}[section]
\newtheorem{proposition}[theorem]{Proposition}
\newtheorem{lemma}[theorem]{Lemma}
\newtheorem{corollary}[theorem]{Corollary}
\newtheorem{condition}[theorem]{Condition}
\theoremstyle{definition}
\newtheorem{definition}[theorem]{Definition}
\newtheorem{remark}[theorem]{Remark}
\newtheorem{example}[theorem]{Example}

\newcommand{\E}{\mathbb{E}}
\newcommand{\Prob}{\mathbb{P}}
\newcommand{\R}{\mathbb{R}}
\newcommand{\1}{\mathbf{1}}
\DeclareMathOperator{\spn}{sp}
\DeclareMathOperator{\supp}{supp}

\title{Uniform $\varepsilon$-Equilibria in Finite Stochastic Games:\\
From Classical Irreducibility to Irreducibility Despite Every Player}
\author{Anonymous draft}
\date{v30}

\begin{document}
\maketitle

\begin{abstract}
This paper has two parts. Part~I records the v28 argument for the classical irreducible regime: if every stationary mixed profile induces an irreducible Markov chain on the original state space $S$, then a stationary average-reward equilibrium yields a uniform $\varepsilon$-equilibrium with explicit $O(1/T)$ regret. The proof is the bias-span argument from the average-reward Markov-decision-process literature: against fixed stationary opponents, each unilateral deviation problem is a finite unichain MDP, the Bellman inequality yields a supermartingale that controls every behavioral deviation, and a stationary average-reward equilibrium gives the compliant lower bound.

Part~II replaces that classical hypothesis by the correct Fudenberg-Yamamoto condition: for each player $i$, the game is \emph{irreducible despite player $i$}. This condition is existential in the opponents' profile and is genuinely weaker than classical irreducibility. Fudenberg and Yamamoto \cite{FudenbergYamamoto2011} show that the opponents may be taken to use statewise uniform randomization, and that under this guardian profile every target state is reached within at most $|S|$ periods with a player-specific probability bounded away from zero, regardless of player $i$'s behavior. We exploit that bounded-return profile by augmenting the state space to
\[
\Omega_K=S\times\{\mathsf R,1,\dots,K\},
\]
where $\mathsf R$ is a return mode and $1,\dots,K$ are work modes. In return mode the profile is forced to play uniform randomization until an anchor state is hit; then the process enters a $K$-stage work block; after stage $K$ it returns to mode $\mathsf R$. Against fixed stationary opponents in the restricted augmented game, both the unrestricted and the anchor-compatible unilateral problems are finite unichain average-reward MDPs on $\Omega_K$. Hence the Part~I Bellman-supermartingale machinery applies on the augmented space.

The remaining comparison is regenerative. A cycle starts at the anchor state in work mode~$1$, contains exactly $K$ work stages, and then a return phase whose expected length is bounded independently of $K$. The unrestricted and restricted MDPs differ only on that return phase. A renewal-ratio comparison yields an $O(1/K)$ gap between the unrestricted and restricted gains, while the bias term on the augmented unichain MDP contributes the usual $O(1/T)$ term. Therefore the restricted equilibrium profile on $\Omega_K$ induces a public finite-memory strategy in the original stochastic game with regret bounded by
\[
O(1/K)+O(H(K)/T),
\]
where $H(K)$ is a bias-span constant on the augmented space and, in fact, $H(K)=O(K)$. Choosing $K$ first and then $T$ proves existence of a uniform $\varepsilon$-equilibrium under the correct Fudenberg-Yamamoto condition. The paper still does not resolve Neyman's conjecture for arbitrary finite stochastic games; it proves the theorem only under this irreducibility-despite-every-player hypothesis.
\end{abstract}

\section{Introduction}

Finite stochastic games, introduced by Shapley \cite{Shapley1953}, are the basic finite-state model of dynamic strategic interaction with publicly observed states. In the zero-sum two-player case, Mertens and Neyman \cite{MertensNeyman1981} proved existence of the uniform value. In the nonzero-sum two-player case, Vieille \cite{Vieille2000a,Vieille2000b} proved existence of uniform equilibrium payoffs. For general multiplayer finite stochastic games, the existence of uniform equilibria remains open; see Solan and Vieille \cite{SolanVieille2015}, Solan \cite{Solan2022}, and the current Game Theory Society open-problem list \cite{OpenProblems2026}.

The purpose of this paper is narrower than that open problem, but broader than v28. Part~I keeps the v28 proof under the classical irreducibility assumption on the original state space $S$. Part~II proves the same uniform-equilibrium conclusion under the \emph{correct} Fudenberg-Yamamoto condition: the game is \emph{irreducible despite every player}. The quantifier pattern matters.

The classical condition used in Part~I is
\[
\forall i\ \forall f_i\ \forall x_{-i}:\ P^{f_i,x_{-i}}\text{ is irreducible on }S.
\]
By contrast, the Fudenberg-Yamamoto condition is existential in the opponents' profile:
\[
\forall i\ \exists x_{-i}^{\mathrm{hub}}\ \forall f_i:\ P^{f_i,x_{-i}^{\mathrm{hub}}}\text{ is irreducible on }S,
\]
and Fudenberg and Yamamoto prove that one may take $x_{-i}^{\mathrm{hub}}$ to be the statewise uniform profile of the opponents \cite[Definition~5 and Lemma~1]{FudenbergYamamoto2011}. Thus the corrected condition is strictly weaker. It does \emph{not} require irreducibility under all stationary profiles; it requires only one guardian profile for the opponents of each player. Example~\ref{ex:strictly-weaker} records a two-state XOR game in which the Fudenberg-Yamamoto condition holds for both players while classical irreducibility fails.

The proof problem under the weaker hypothesis is exactly the multichain obstruction that v28 avoided. Without classical irreducibility, a unilateral stationary deviation can in principle create several recurrent classes with different gains, and then a single scalar gain on the original state space need not control all deviations. The Part~II repair is to move the bias argument to a \emph{return-to-anchor augmentation}. The augmented process forces the game to spend short episodes in a guardian mode where all players use uniform randomization. Those episodes recur, so they cannot be treated as a one-time transient prefix. Instead, the regret bound has the form
\[
O(1/K)+O(H(K)/T),
\]
where $K$ is the work-block length and $H(K)$ is the bias-span term on the augmented state space.

The key idea is simple. Fix an anchor state $s^\bullet\in S$ and an integer $K\ge1$. Add a public mode variable with values
\[
\mathsf R,1,2,\dots,K.
\]
Mode $\mathsf R$ is the return mode. There all players are forced to use the uniform profile. The mode stays at $\mathsf R$ until the physical state hits $s^\bullet$, and then the process enters work mode~$1$. In work modes $1,\dots,K$ players use stationary mixed actions on the augmented state space. After work mode $K$, the process returns to $\mathsf R$. This construction turns the original game into a finite stochastic game on
\[
\Omega_K=S\times\{\mathsf R,1,\dots,K\}.
\]

The Fudenberg-Yamamoto bounded-return lemma implies that from every augmented state there is a uniformly positive probability of reaching the gate state $(s^\bullet,1)$ within a bounded number of stages, regardless of the deviator's behavior. Hence every recurrent class of the unilateral control problem on $\Omega_K$ must contain $(s^\bullet,1)$. The unrestricted and restricted unilateral problems are therefore unichain MDPs on $\Omega_K$. This is the decisive point: once the augmented chain is unichain, transient states are handled automatically by the bias on $\Omega_K$, exactly as in Part~I.

We prove two theorems. The first is the v28 statement, kept as Part~I.

\begin{theorem}[Classical irreducible case]\label{thm:intro-main}
Let $G$ be a finite stochastic game satisfying the classical irreducibility condition of Part~I. Then for every $\varepsilon>0$ there exist a stationary mixed strategy profile $x^*$ and an integer $T_0$ such that for every initial state $s\in S$, every horizon $T\ge T_0$, every player $i$, and every behavioral deviation $\tau_i$ of player $i$,
\[
\gamma_i^T(s,x^*)\ge \gamma_i^T\bigl(s,(\tau_i,x_{-i}^*)\bigr)-\varepsilon.
\]
More precisely,
\[
\gamma_i^T\bigl(s,(\tau_i,x_{-i}^*)\bigr)-\gamma_i^T(s,x^*)
\le \frac{2\max_j\spn(h_j^*)}{T}
\]
for all $T\ge1$, where $h_j^*$ is a normalized bias for player $j$ against $x_{-j}^*$.
\end{theorem}

The main new statement is the Part~II theorem.

\begin{theorem}[Fudenberg-Yamamoto case]\label{thm:intro-fy}
Assume the Fudenberg-Yamamoto condition of Part~II. Fix an anchor state $s^\bullet\in S$, let $p_i>0$ be the bounded-return constants supplied by Proposition~\ref{prop:guardian}, and set
\[
M_i:=\frac{|S|}{p_i},
\qquad
M:=\max_{i\in N}M_i.
\]
Then for every integer $K\ge1$ there exists a public finite-memory strategy profile $\sigma^K$ in the original stochastic game such that for every horizon $T\ge1$,
\[
\sup_{i\in N}\sup_{s\in S}\sup_{\tau_i}
\Bigl[\gamma_i^T\bigl(s,(\tau_i,\sigma_{-i}^K)\bigr)-\gamma_i^T(s,\sigma^K)\Bigr]
\le \frac{2M}{K}+\frac{H(K)}{T},
\]
where $H(K)<\infty$ and, by Proposition~\ref{prop:aug-bias-span}, $H(K)=O(K)$.

Consequently, for every $\varepsilon>0$, if
\[
K_\varepsilon:=\left\lceil \frac{4M}{\varepsilon}\right\rceil
\qquad\text{and}\qquad
T_0(\varepsilon):=\left\lceil \frac{2H(K_\varepsilon)}{\varepsilon}\right\rceil,
\]
then $\sigma^{K_\varepsilon}$ is a uniform $\varepsilon$-equilibrium.
\end{theorem}

The difference between the two parts is therefore structural, not philosophical. In Part~I the original state space already yields a unichain unilateral problem. In Part~II one must create a unichain problem by augmenting the state space and by forcing recurrent returns to the guardian mode. Once that is done, the same Bellman supermartingale controls arbitrary behavioral deviations.

The paper is organized as follows. Part~I, namely Sections~\ref{sec:model}--\ref{sec:proof-main}, records the classical irreducible-case proof from v28. Part~II begins in Section~\ref{sec:FY} with the correct Fudenberg-Yamamoto condition and its bounded-return guardian profile. Section~\ref{sec:augmentation} defines the return-to-anchor augmentation. Section~\ref{sec:aug-metagame} proves existence of a restricted stationary average-reward equilibrium on the augmented space. Section~\ref{sec:regenerative} contains the regenerative $O(1/K)$ comparison and the bias-span bound. Section~\ref{sec:proof-fy} proves Theorem~\ref{thm:intro-fy}. Section~\ref{sec:discussion-new} closes with scope and remaining open problems.

\part{Classical irreducibility}

\section{Finite stochastic games and the classical irreducibility condition}\label{sec:model}

\subsection{Game model}

A finite stochastic game is given by
\[
G=(N,S,(A_i)_{i\in N},u,P),
\]
where:
\begin{itemize}[leftmargin=2.3em]
\item $N=\{1,\dots,n\}$ is the finite set of players,
\item $S$ is a finite nonempty set of states,
\item for each player $i$ and each state $s\in S$, the action set $A_i(s)$ is finite and nonempty,
\item for each player $i$, the stage payoff function satisfies
\[
 u_i(s,a)\in[0,1]\qquad\text{for all }s\in S,\ a=(a_j)_{j\in N}\in A(s):=\prod_{j\in N}A_j(s),
\]
\item $P(\cdot\mid s,a)\in\Delta(S)$ is the transition law from state $s$ under action profile $a$.
\end{itemize}

At stage $t$, the current state $s_t$ is publicly observed, the players simultaneously choose an action profile $a_t=(a_{1,t},\dots,a_{n,t})\in A(s_t)$, players receive payoffs $u_i(s_t,a_t)$, and the next state is drawn according to $P(\cdot\mid s_t,a_t)$.

A \emph{behavioral strategy} of player $i$ is a rule $\tau_i$ assigning to every finite public history
\[
h_t=(s_1,a_1,s_2,a_2,\dots,s_t)
\]
a probability distribution $\tau_i(\cdot\mid h_t)\in\Delta(A_i(s_t))$. A strategy profile is denoted by $\sigma=(\sigma_1,\dots,\sigma_n)$.

For stagewise arguments we use the \emph{pre-action public filtration}
\[
\mathcal F_t:=\sigma(s_1,a_1,\dots,s_{t-1},a_{t-1},s_t),\qquad t\ge1.
\]
Thus $\mathcal F_t$ is exactly the publicly observed history up to stage $t$, just before the stage-$t$ action is chosen. In particular, $s_t$ and the mixed action $\tau_i(\cdot\mid h_t)$ prescribed by any behavioral strategy are $\mathcal F_t$-measurable.

For $T\ge1$, the expected $T$-stage average payoff of player $i$ from initial state $s$ under profile $\sigma$ is
\[
\gamma_i^T(s,\sigma)
=\E_s^{\sigma}\!\left[\frac1T\sum_{t=1}^{T}u_i(s_t,a_t)\right].
\]

\begin{definition}
A strategy profile $\sigma$ is a \emph{uniform $\varepsilon$-equilibrium} if there exists $T_0\in\mathbb N$ such that for every horizon $T\ge T_0$, every initial state $s\in S$, every player $i$, and every behavioral deviation $\tau_i$ of player $i$,
\[
\gamma_i^T(s,\sigma)\ge\gamma_i^T\bigl(s,(\tau_i,\sigma_{-i})\bigr)-\varepsilon.
\]
\end{definition}

We will use the following stationary strategy spaces.
\[
X_i:=\prod_{s\in S}\Delta(A_i(s)),\qquad X:=\prod_{i\in N}X_i.
\]
An element $x_i\in X_i$ is a \emph{stationary mixed strategy}. We write $x_i(a_i\mid s)$ for the probability assigned to action $a_i\in A_i(s)$. A \emph{pure stationary strategy} of player $i$ is a map $f_i:S\to\bigcup_{s}A_i(s)$ such that $f_i(s)\in A_i(s)$ for all $s$; the finite set of such strategies is denoted by
\[
F_i:=\prod_{s\in S}A_i(s).
\]

Given $x=(x_1,\dots,x_n)\in X$, define its stagewise product distribution by
\[
x(a\mid s):=\prod_{j\in N}x_j(a_j\mid s),\qquad a\in A(s),\ s\in S.
\]
Then the induced transition matrix on $S$ is
\[
P^x(s,s'):=\sum_{a\in A(s)}x(a\mid s)P(s'\mid s,a),\qquad s,s'\in S.
\]
For a pure stationary strategy $f_i\in F_i$ and an opponents' stationary mixed profile $x_{-i}\in X_{-i}$, we similarly write
\[
P^{f_i,x_{-i}}(s,s'):=\sum_{a_{-i}\in A_{-i}(s)}x_{-i}(a_{-i}\mid s)
P\bigl(s'\mid s,f_i(s),a_{-i}\bigr).
\]

\subsection{Condition A and classical irreducibility}

\begin{condition}[Classical irreducibility]\label{cond:A}
For every player $i\in N$, every pure stationary strategy $f_i\in F_i$, and every stationary mixed profile $x_{-i}\in X_{-i}$ of the other players, the Markov chain on $S$ with transition matrix $P^{f_i,x_{-i}}$ is irreducible.
\end{condition}

\begin{remark}
Condition~\ref{cond:A} is imposed on the original game and on the original state space $S$. A simple sufficient condition is a Doeblin-type lower bound: if there exist $\alpha>0$ and a probability measure $q$ with full support on $S$ such that
\[
P(\cdot\mid s,a)\ge \alpha q(\cdot)\qquad\text{for every }s\in S\text{ and }a\in A(s),
\]
then Condition~\ref{cond:A} holds automatically. The point of Condition~\ref{cond:A}, however, is that it does \emph{not} require such a uniform minorization. It only requires irreducibility of the full chain on $S$ under every unilateral pure stationary deviation against stationary mixed opponents.
\end{remark}

The first basic consequence is that mixed stationary strategies do not destroy irreducibility.

\begin{proposition}\label{prop:mixed-irreducible}
Assume Condition~\ref{cond:A}. Then for every stationary mixed profile $x\in X$, the Markov chain on $S$ with transition matrix $P^x$ is irreducible.
\end{proposition}

\begin{proof}
Fix $x\in X$ and fix one player $i\in N$. For each pure stationary strategy $f_i\in F_i$, define
\[
\lambda_{f_i}(x_i):=\prod_{s\in S}x_i\bigl(f_i(s)\mid s\bigr).
\]
Because the action choices at different states factor, we have
\[
\sum_{f_i\in F_i}\lambda_{f_i}(x_i)=1.
\]
Moreover, a direct expansion over pure stationary strategies gives the matrix identity
\begin{equation}\label{eq:convex-combination}
P^x=\sum_{f_i\in F_i}\lambda_{f_i}(x_i)\,P^{f_i,x_{-i}}.
\end{equation}
Indeed, for each pair of states $s,s'$, the $(s,s')$-entry of the right-hand side equals
\begin{align*}
\sum_{f_i\in F_i}\lambda_{f_i}(x_i)P^{f_i,x_{-i}}(s,s')
&=\sum_{f_i\in F_i}\left(\prod_{z\in S}x_i(f_i(z)\mid z)\right)
\sum_{a_{-i}}x_{-i}(a_{-i}\mid s)P(s'\mid s,f_i(s),a_{-i})\\
&=\sum_{a_i\in A_i(s)}x_i(a_i\mid s)\sum_{a_{-i}}x_{-i}(a_{-i}\mid s)P(s'\mid s,a_i,a_{-i})\\
&=P^x(s,s').
\end{align*}
Choose $f_i^\sharp\in F_i$ such that $x_i(f_i^\sharp(s)\mid s)>0$ for every state $s$. Then $\lambda_{f_i^\sharp}(x_i)>0$. By Condition~\ref{cond:A}, the chain $P^{f_i^\sharp,x_{-i}}$ is irreducible. Therefore, for every pair of states $p,q\in S$, there exists $m\ge1$ such that
\[
\bigl(P^{f_i^\sharp,x_{-i}}\bigr)^m(p,q)>0.
\]
From \eqref{eq:convex-combination} we have the entrywise inequality
\[
P^x\ge \lambda_{f_i^\sharp}(x_i)\,P^{f_i^\sharp,x_{-i}},
\]
whence, by monotonicity of matrix multiplication for nonnegative matrices,
\[
(P^x)^m\ge \lambda_{f_i^\sharp}(x_i)^m\bigl(P^{f_i^\sharp,x_{-i}}\bigr)^m.
\]
In particular,
\[
(P^x)^m(p,q)\ge \lambda_{f_i^\sharp}(x_i)^m\bigl(P^{f_i^\sharp,x_{-i}}\bigr)^m(p,q)>0.
\]
Thus $P^x$ is irreducible.
\end{proof}

\begin{remark}[Condition~\ref{cond:A} is the classical irreducibility assumption]\label{rem:classical-irreducibility}
Proposition~\ref{prop:mixed-irreducible} shows that Condition~\ref{cond:A} implies irreducibility under every stationary mixed profile $x\in X$. The converse is immediate, because a unilateral pure stationary strategy against stationary mixed opponents is a special case of a stationary mixed profile. Thus Condition~\ref{cond:A} is exactly the classical irreducibility requirement for stationary profiles. We retain the unilateral formulation because it interfaces directly with the MDP arguments below. Part~II will introduce a genuinely weaker existential condition of Fudenberg and Yamamoto.
\end{remark}

For the unilateral optimization problems that will appear later, this immediately implies a unichain property.

\begin{corollary}\label{cor:unichain-from-A}
Assume Condition~\ref{cond:A}, fix a player $i$, and fix an opponents' stationary mixed profile $x_{-i}\in X_{-i}$. Consider the MDP faced by player $i$ on state space $S$, with action set $A_i(s)$ at state $s$, transition kernel
\[
Q_i^{x_{-i}}(s'\mid s,a_i):=\sum_{a_{-i}\in A_{-i}(s)}x_{-i}(a_{-i}\mid s)P(s'\mid s,a_i,a_{-i}),
\]
and reward
\[
r_i^{x_{-i}}(s,a_i):=\sum_{a_{-i}\in A_{-i}(s)}x_{-i}(a_{-i}\mid s)u_i(s,a_i,a_{-i}).
\]
Then every pure stationary policy of this MDP induces an irreducible Markov chain on $S$. In particular, the MDP is unichain; indeed it is irreducible under every pure stationary policy.
\end{corollary}

\begin{proof}
A pure stationary policy of the MDP is exactly a pure stationary strategy $f_i\in F_i$ of player $i$ in the stochastic game. Its induced transition matrix is $P^{f_i,x_{-i}}$, which is irreducible by Condition~\ref{cond:A}.
\end{proof}

\section{Average-reward MDP facts under fixed opponents}\label{sec:mdp}

Fix a player $i\in N$ and an opponents' stationary mixed profile $x_{-i}\in X_{-i}$. The controlled process of player $i$ is the finite MDP from Corollary~\ref{cor:unichain-from-A}.

For a stationary mixed policy $\mu_i\in X_i$, its induced transition matrix is
\[
P^{\mu_i,x_{-i}}(s,s'):=\sum_{a_i\in A_i(s)}\mu_i(a_i\mid s)
Q_i^{x_{-i}}(s'\mid s,a_i),
\]
and its expected one-step reward in state $s$ is
\[
\bar r_i^{\mu_i,x_{-i}}(s):=\sum_{a_i\in A_i(s)}\mu_i(a_i\mid s)r_i^{x_{-i}}(s,a_i).
\]
By Proposition~\ref{prop:mixed-irreducible}, the chain $P^{\mu_i,x_{-i}}$ is irreducible, so it has a unique invariant distribution $\pi^{\mu_i,x_{-i}}$ and an average reward
\[
g_i(\mu_i,x_{-i})
:=\sum_{s\in S}\pi^{\mu_i,x_{-i}}(s)\,\bar r_i^{\mu_i,x_{-i}}(s),
\]
independent of the initial state.

\subsection{Gain and bias}

Fix a reference state $s^\dagger\in S$ once and for all. The following is standard finite-state average-reward MDP theory; see Puterman \cite[Chapter~8]{Puterman1994} or Neyman \cite[Section~3]{Neyman2003MarkovChains}.

\begin{theorem}[Average-reward optimality equation]\label{thm:aroe}
Assume Condition~\ref{cond:A}. For each player $i$ and each opponents' stationary mixed profile $x_{-i}\in X_{-i}$, there exist a scalar $g_i^*(x_{-i})\in\R$ and a function $h_i^{x_{-i}}:S\to\R$, normalized by
\[
h_i^{x_{-i}}(s^\dagger)=0,
\]
such that
\begin{equation}\label{eq:bellman-equality}
g_i^*(x_{-i})+h_i^{x_{-i}}(s)
=\max_{a_i\in A_i(s)}\Bigl[r_i^{x_{-i}}(s,a_i)+\sum_{s'\in S}Q_i^{x_{-i}}(s'\mid s,a_i)h_i^{x_{-i}}(s')\Bigr]
\end{equation}
for every $s\in S$.

The number $g_i^*(x_{-i})$ is the optimal average reward of player $i$ against $x_{-i}$ over all admissible policies, including history-dependent randomized policies. Moreover, every pure stationary selector of the maximizers in \eqref{eq:bellman-equality} is average-reward optimal. Since $S$ is finite, every such normalized bias has finite span.
\end{theorem}

\begin{proof}
By Corollary~\ref{cor:unichain-from-A}, the induced MDP is finite and unichain. The existence of a normalized solution $(g_i^*(x_{-i}),h_i^{x_{-i}})$ to the average-reward optimality equation, together with the optimality of stationary maximizing selectors, is standard finite-state average-reward MDP theory; see Puterman \cite[Chapter~8]{Puterman1994} or Neyman \cite[Section~3]{Neyman2003MarkovChains}.
\end{proof}

Write
\[
A_i^*(s;x_{-i})
:=\arg\max_{a_i\in A_i(s)}\Bigl[r_i^{x_{-i}}(s,a_i)+\sum_{s'\in S}Q_i^{x_{-i}}(s'\mid s,a_i)h_i^{x_{-i}}(s')\Bigr].
\]
Each set $A_i^*(s;x_{-i})$ is nonempty by finiteness.

\subsection{The supermartingale bound for arbitrary deviations}

The next proposition is the key estimate in the entire argument. It bounds the finite-horizon payoff of \emph{every} unilateral deviation, not merely stationary deviations.

\begin{proposition}[Supermartingale upper bound]\label{prop:supermartingale}
Assume Condition~\ref{cond:A}. Fix a player $i$ and a stationary mixed opponents' profile $x_{-i}\in X_{-i}$. Let $g_i^*=g_i^*(x_{-i})$ and $h_i=h_i^{x_{-i}}$ be the gain and normalized bias from Theorem~\ref{thm:aroe}. Then for every initial state $s\in S$, every horizon $T\ge1$, and every behavioral strategy $\tau_i$ of player $i$,
\begin{equation}\label{eq:deviation-upper}
\gamma_i^T\bigl(s,(\tau_i,x_{-i})\bigr)
\le g_i^*+\frac{\spn(h_i)}{T}.
\end{equation}
Equivalently, if $u_t:=u_i(s_t,a_t)$ and
\[
M_t:=h_i(s_t)+\sum_{k=1}^{t-1}(g_i^*-u_k),\qquad t\ge1,
\]
then $(M_t)_{t\ge1}$ is a supermartingale under $\Prob_s^{(\tau_i,x_{-i})}$.
\end{proposition}

\begin{proof}
Fix $s$, $T$, and $\tau_i$. Let $(\mathcal F_t)_{t\ge1}$ be the pre-action public filtration from Section~\ref{sec:model}, so
\[
\mathcal F_t=\sigma(s_1,a_1,\dots,s_{t-1},a_{t-1},s_t).
\]
Write
\[
\mu_t(\cdot):=\tau_i(\cdot\mid h_t),
\]
where $h_t=(s_1,a_1,\dots,s_t)$. Then $s_t$ and $\mu_t$ are $\mathcal F_t$-measurable. Under the profile $(\tau_i,x_{-i})$, conditional on $\mathcal F_t$, player $i$'s stage-$t$ action is sampled from $\mu_t(\cdot)$, the opponents' action profile is sampled from $x_{-i}(\cdot\mid s_t)$, and then $s_{t+1}$ is drawn from $P(\cdot\mid s_t,a_t)$.

By Theorem~\ref{thm:aroe}, for every state $s$ and every action $a_i\in A_i(s)$,
\[
r_i^{x_{-i}}(s,a_i)+\sum_{s'}Q_i^{x_{-i}}(s'\mid s,a_i)h_i(s')\le g_i^*+h_i(s).
\]
Applying this inequality at the random state $s_t$ and averaging over the $\mathcal F_t$-measurable mixed action $\mu_t$ yields
\begin{align*}
\E\bigl[u_t+h_i(s_{t+1})\mid\mathcal F_t\bigr]
&=\sum_{a_i\in A_i(s_t)}\mu_t(a_i)\left(r_i^{x_{-i}}(s_t,a_i)+\sum_{s'}Q_i^{x_{-i}}(s'\mid s_t,a_i)h_i(s')\right)\\
&\le g_i^*+h_i(s_t).
\end{align*}
Rearranging gives
\[
\E[M_{t+1}\mid \mathcal F_t]\le M_t,
\]
so $(M_t)$ is a supermartingale.

Taking expectations and using $M_1=h_i(s)$, we obtain
\[
\E[h_i(s_{T+1})]+Tg_i^*-\E\Bigl[\sum_{t=1}^{T}u_t\Bigr]\le h_i(s).
\]
Hence
\[
\E\Bigl[\sum_{t=1}^{T}u_t\Bigr]\le Tg_i^*+h_i(s)-\E[h_i(s_{T+1})].
\]
Because the difference between any two values of $h_i$ is bounded by its span,
\[
h_i(s)-\E[h_i(s_{T+1})]\le \max_{z\in S}h_i(z)-\min_{z\in S}h_i(z)=\spn(h_i).
\]
Dividing by $T$ yields \eqref{eq:deviation-upper}.
\end{proof}

\subsection{Optimal stationary responses and the exact finite-horizon identity}

The next result identifies the stationary optimal responses to $x_{-i}$ and gives the matching lower bound along such responses.

\begin{proposition}[Optimal stationary responses]\label{prop:optimal-stationary}
Assume Condition~\ref{cond:A}. Fix a player $i$ and a stationary mixed opponents' profile $x_{-i}\in X_{-i}$. Let $g_i^*=g_i^*(x_{-i})$ and $h_i=h_i^{x_{-i}}$ be as in Theorem~\ref{thm:aroe}.

\begin{enumerate}[label=\textup{(\alph*)},leftmargin=2.4em]
\item If a stationary mixed strategy $\mu_i\in X_i$ satisfies
\[
\supp \mu_i(\cdot\mid s)\subseteq A_i^*(s;x_{-i})\qquad\text{for every }s\in S,
\]
then for every initial state $s\in S$ and every horizon $T\ge1$,
\begin{equation}\label{eq:exact-identity}
\gamma_i^T\bigl(s,(\mu_i,x_{-i})\bigr)
= g_i^*+\frac{h_i(s)-\E_s^{(\mu_i,x_{-i})}[h_i(s_{T+1})]}{T}.
\end{equation}
In particular,
\begin{equation}\label{eq:exact-lower-bound}
\left|\gamma_i^T\bigl(s,(\mu_i,x_{-i})\bigr)-g_i^*\right|\le \frac{\spn(h_i)}{T},
\end{equation}
so $\mu_i$ is average-reward optimal.

\item Conversely, if a stationary mixed strategy $\mu_i\in X_i$ is average-reward optimal against $x_{-i}$, then
\[
\supp \mu_i(\cdot\mid s)\subseteq A_i^*(s;x_{-i})\qquad\text{for every }s\in S.
\]
Hence the stationary optimal responses are exactly the stationary mixed strategies whose support is contained statewise in the maximizing action sets.
\end{enumerate}
\end{proposition}

\begin{proof}
\textit{Part (a).} Let $\mu_i$ satisfy the support condition. Averaging the Bellman equality \eqref{eq:bellman-equality} over $\mu_i(\cdot\mid s)$ gives, for each state $s$,
\[
g_i^*+h_i(s)
=\bar r_i^{\mu_i,x_{-i}}(s)+\sum_{s'\in S}P^{\mu_i,x_{-i}}(s,s')h_i(s').
\]
Therefore, under the profile $(\mu_i,x_{-i})$,
\[
\E\bigl[u_t+h_i(s_{t+1})\mid s_t\bigr]=g_i^*+h_i(s_t).
\]
Taking full expectations and summing from $t=1$ to $T$ yields
\[
\E\Bigl[\sum_{t=1}^{T}u_t\Bigr]+\E[h_i(s_{T+1})]=Tg_i^*+h_i(s),
\]
which is exactly \eqref{eq:exact-identity}. The bound \eqref{eq:exact-lower-bound} follows immediately because the numerator is bounded in absolute value by $\spn(h_i)$. Thus $\mu_i$ attains the optimal average reward $g_i^*$.

\smallskip
\noindent
\textit{Part (b).} Let $\mu_i$ be stationary and average-reward optimal against $x_{-i}$. Define the nonnegative slack function
\[
d(s,a_i):=g_i^*+h_i(s)-r_i^{x_{-i}}(s,a_i)-\sum_{s'\in S}Q_i^{x_{-i}}(s'\mid s,a_i)h_i(s').
\]
By construction, $d(s,a_i)\ge0$ for every state-action pair. Let $\pi$ be the unique invariant distribution of the irreducible chain $P^{\mu_i,x_{-i}}$. Then
\begin{align*}
\sum_{s\in S}\pi(s)\sum_{a_i\in A_i(s)}\mu_i(a_i\mid s)d(s,a_i)
&=\sum_{s}\pi(s)\left(g_i^*+h_i(s)-\bar r_i^{\mu_i,x_{-i}}(s)-\sum_{s'}P^{\mu_i,x_{-i}}(s,s')h_i(s')\right)\\
&=g_i^*-\sum_s\pi(s)\bar r_i^{\mu_i,x_{-i}}(s)\\
&=g_i^*-g_i(\mu_i,x_{-i})\\
&=0.
\end{align*}
The second equality uses the invariance identity $\sum_s\pi(s)h_i(s)=\sum_{s,s'}\pi(s)P^{\mu_i,x_{-i}}(s,s')h_i(s')$, and the last equality uses optimality. Since all terms are nonnegative and, because the chain is irreducible, $\pi(s)>0$ for every state $s$, we must have
\[
\sum_{a_i}\mu_i(a_i\mid s)d(s,a_i)=0\qquad\text{for every }s\in S.
\]
Thus whenever $\mu_i(a_i\mid s)>0$, necessarily $d(s,a_i)=0$, which means $a_i\in A_i^*(s;x_{-i})$. This proves the support inclusion.
\end{proof}

\section{The average-reward meta-game}\label{sec:metagame}

\subsection{Definition and continuity of the payoff map}

For a stationary mixed profile $x\in X$, define the expected one-step payoff of player $i$ in state $s$ by
\[
\bar u_i^x(s):=\sum_{a\in A(s)}x(a\mid s)u_i(s,a).
\]
By Proposition~\ref{prop:mixed-irreducible}, the chain $P^x$ is irreducible. Let $\pi^x$ denote its unique invariant distribution. We define the average-reward meta-game payoff of player $i$ at $x$ by
\begin{equation}\label{eq:meta-payoff}
g_i(x):=\sum_{s\in S}\pi^x(s)\bar u_i^x(s).
\end{equation}
Because $P^x$ is irreducible, this quantity is equal to the unique average reward of player $i$ under the stationary profile $x$, independent of the initial state.

\begin{proposition}[Continuity of invariant distributions and average rewards]\label{prop:continuity}
Assume Condition~\ref{cond:A}. For every player $i$, the map
\[
x\longmapsto \pi^x\in \Delta(S)
\]
is continuous on $X$, and therefore the payoff map
\[
x\longmapsto g_i(x)
\]
is continuous on $X$.
\end{proposition}

\begin{proof}
The map $x\mapsto P^x$ is multilinear, hence continuous. Let $(x^m)_{m\ge1}$ be a sequence in $X$ converging to $x\in X$. Write $\pi^m:=\pi^{x^m}$. Since $\Delta(S)$ is compact, every subsequence of $(\pi^m)$ admits a further convergent subsequence. Let $\pi^{m_k}\to \bar\pi$.

For each $k$, the invariance identity is
\[
\pi^{m_k}P^{x^{m_k}}=\pi^{m_k}.
\]
Passing to the limit and using continuity of matrix multiplication gives
\[
\bar\pi P^x=\bar\pi.
\]
Thus $\bar\pi$ is an invariant distribution for $P^x$. By Proposition~\ref{prop:mixed-irreducible}, $P^x$ is irreducible, so its invariant distribution is unique. Hence $\bar\pi=\pi^x$.

We have shown that every convergent subsequence of $(\pi^m)$ converges to $\pi^x$. Since $\Delta(S)$ is compact, this implies the whole sequence converges: $\pi^m\to \pi^x$. Therefore $x\mapsto \pi^x$ is continuous.

The map $x\mapsto \bar u_i^x(s)$ is continuous for each state $s$ because it is multilinear in the stationary mixed strategies. Formula \eqref{eq:meta-payoff} then shows that $g_i$ is continuous as the composition of continuous operations.
\end{proof}

The next proposition records the ergodic interpretation of $g_i(x)$. The proof is standard but useful because it links the meta-game payoff to the actual long-run behavior of play.

\begin{proposition}[Ergodic theorem for stationary profiles]\label{prop:ergodic}
Assume Condition~\ref{cond:A}. Fix a stationary mixed profile $x\in X$, an initial state $s\in S$, and a player $i$. Then under $\Prob_s^x$,
\[
\frac1T\sum_{t=1}^{T}u_i(s_t,a_t)\longrightarrow g_i(x)
\qquad\text{almost surely and in }L^1.
\]
\end{proposition}

\begin{proof}
Because $P^x$ is a finite irreducible Markov chain, the ergodic theorem for finite irreducible chains implies that for every bounded function $f:S\to\R$,
\[
\frac1T\sum_{t=1}^{T}f(s_t)\longrightarrow\sum_{z\in S}\pi^x(z)f(z)
\qquad\text{almost surely and in }L^1.
\]
Apply this to $f=\bar u_i^x$. Then
\begin{equation}\label{eq:ergodic-expected-reward}
\frac1T\sum_{t=1}^{T}\bar u_i^x(s_t)\longrightarrow \sum_{z\in S}\pi^x(z)\bar u_i^x(z)=g_i(x)
\qquad\text{almost surely and in }L^1.
\end{equation}

Next define
\[
d_t:=u_i(s_t,a_t)-\bar u_i^x(s_t),\qquad t\ge1.
\]
Let
\[
\mathcal G_t:=\sigma(s_1,a_1,s_2,a_2,\dots,s_t,a_t,s_{t+1}),\qquad t\ge0,
\]
with the convention that $\mathcal G_0=\sigma(s_1)$. Then $d_t$ is $\mathcal G_t$-measurable and, because under the stationary profile $x$ the action at stage $t$ is drawn according to $x(\cdot\mid s_t)$,
\[
\E[d_t\mid \mathcal G_{t-1}]=0.
\]
Hence the partial sums $D_T:=\sum_{t=1}^{T}d_t$ form a martingale with bounded increments. By the strong law for bounded martingale differences,
\[
\frac{D_T}{T}\longrightarrow 0
\qquad\text{almost surely and in }L^1.
\]
Combining this with \eqref{eq:ergodic-expected-reward} yields
\[
\frac1T\sum_{t=1}^{T}u_i(s_t,a_t)
=\frac1T\sum_{t=1}^{T}\bar u_i^x(s_t)+\frac{D_T}{T}
\longrightarrow g_i(x)
\]
almost surely and in $L^1$.
\end{proof}

\subsection{Best responses in the meta-game}

Fix a player $i$. The best-response correspondence of player $i$ in the average-reward meta-game is
\[
B_i(x_{-i}):=\arg\max_{\mu_i\in X_i} g_i(\mu_i,x_{-i}),\qquad x_{-i}\in X_{-i}.
\]

\begin{proposition}[Structure of best responses]\label{prop:best-response-structure}
Assume Condition~\ref{cond:A}. Fix a player $i$ and $x_{-i}\in X_{-i}$. Then
\begin{equation}\label{eq:best-response-product}
B_i(x_{-i})=
\Bigl\{\mu_i\in X_i:\supp\mu_i(\cdot\mid s)\subseteq A_i^*(s;x_{-i})\text{ for every }s\in S\Bigr\}.
\end{equation}
In particular, $B_i(x_{-i})$ is nonempty, compact, and convex.
\end{proposition}

\begin{proof}
Because $X_i$ is compact and $\mu_i\mapsto g_i(\mu_i,x_{-i})$ is continuous by Proposition~\ref{prop:continuity}, the set $B_i(x_{-i})$ is nonempty and compact.

By Proposition~\ref{prop:optimal-stationary}, a stationary mixed strategy $\mu_i$ is average-reward optimal against $x_{-i}$ if and only if its support is contained statewise in the maximizing sets $A_i^*(s;x_{-i})$. This proves \eqref{eq:best-response-product}. The right-hand side is a product of simplices
\[
\prod_{s\in S}\Delta\bigl(A_i^*(s;x_{-i})\bigr),
\]
hence convex and compact.
\end{proof}

\begin{proposition}[Closed graph of the best-response correspondence]\label{prop:best-response-closed-graph}
Assume Condition~\ref{cond:A}. For each player $i$, the correspondence $B_i:X_{-i}\rightrightarrows X_i$ has closed graph. Consequently it is upper hemicontinuous.
\end{proposition}

\begin{proof}
Let $x_{-i}^m\to x_{-i}$ in $X_{-i}$, let $\mu_i^m\to \mu_i$ in $X_i$, and suppose $\mu_i^m\in B_i(x_{-i}^m)$ for every $m$. Then for every fixed $\eta_i\in X_i$,
\[
g_i(\mu_i^m,x_{-i}^m)\ge g_i(\eta_i,x_{-i}^m).
\]
Passing to the limit and using continuity of $g_i$ from Proposition~\ref{prop:continuity}, we get
\[
g_i(\mu_i,x_{-i})\ge g_i(\eta_i,x_{-i})\qquad\text{for every }\eta_i\in X_i.
\]
Hence $\mu_i\in B_i(x_{-i})$, so the graph is closed.

Because the domain $X_{-i}$ is compact and the values of $B_i$ are compact by Proposition~\ref{prop:best-response-structure}, closed graph implies upper hemicontinuity.
\end{proof}

\begin{remark}[Occupation-measure formulation]\label{rem:occupation}
For fixed $i$ and $x_{-i}$, the average-reward MDP can also be written as a linear program over state-action occupation measures. One maximizes
\[
\sum_{s\in S}\sum_{a_i\in A_i(s)}m(s,a_i)\,r_i^{x_{-i}}(s,a_i)
\]
subject to the balance equations
\[
\sum_{a_i}m(s,a_i)=\sum_{z\in S}\sum_{b_i\in A_i(z)}m(z,b_i)Q_i^{x_{-i}}(s\mid z,b_i),
\qquad s\in S,
\]
together with $m\ge0$ and $\sum_{s,a_i}m(s,a_i)=1$. The feasible set is convex, so the set of $\eta$-optimal \emph{occupation measures} is convex for every $\eta\ge0$. We will not use this formulation below, because Proposition~\ref{prop:best-response-structure} already gives the stronger policy-space description of exact best responses that is needed for the fixed-point argument.
\end{remark}

\subsection{Existence of a stationary average-reward equilibrium}

Stationary average-reward equilibrium existence in this irreducible setting is classical; see Rogers \cite{Rogers1969}, Sobel \cite{Sobel1971}, and Federgruen \cite{Federgruen1978}. We record the short fixed-point argument here only for completeness and because it matches the notation used later in the $O(1/T)$ regret estimate.

We now apply a standard fixed-point theorem to the product best-response correspondence.

\begin{theorem}[Kakutani-Fan-Glicksberg fixed-point theorem]\label{thm:kfg}
Let $K$ be a nonempty compact convex subset of a Euclidean space, and let $\Phi:K\rightrightarrows K$ be a correspondence with nonempty compact convex values and closed graph. Then there exists $x\in K$ such that $x\in \Phi(x)$.
\end{theorem}

\begin{proof}
This is the classical Kakutani-Fan-Glicksberg fixed-point theorem; see Glicksberg \cite{Glicksberg1952}.
\end{proof}

\begin{theorem}[Stationary equilibrium of the average-reward meta-game]\label{thm:meta-equilibrium}
Assume Condition~\ref{cond:A}. Then there exists a stationary mixed profile $x^*\in X$ such that for every player $i$,
\[
g_i(x^*)\ge g_i(\mu_i,x_{-i}^*)\qquad\text{for every }\mu_i\in X_i.
\]
Equivalently, $x_i^*\in B_i(x_{-i}^*)$ for every player $i$.
\end{theorem}

\begin{proof}
Define the product correspondence $B:X\rightrightarrows X$ by
\[
B(x):=\prod_{i\in N}B_i(x_{-i}).
\]
The set $X$ is a nonempty compact convex subset of a finite-dimensional Euclidean space. By Propositions~\ref{prop:best-response-structure} and \ref{prop:best-response-closed-graph}, each $B_i$ has nonempty compact convex values and closed graph. Therefore $B$ itself has nonempty compact convex values and closed graph.

By Theorem~\ref{thm:kfg}, there exists $x^*\in X$ with $x^*\in B(x^*)$. This means precisely that $x_i^*\in B_i(x_{-i}^*)$ for every player $i$, or equivalently that each $x_i^*$ maximizes $g_i(\mu_i,x_{-i}^*)$ over $\mu_i\in X_i$.
\end{proof}

\section{Proof of the main theorem}\label{sec:proof-main}

We can now prove the uniform-equilibrium theorem announced in the introduction. The existence of the stationary average-reward equilibrium used in the first line is classical in the irreducible regime and was also rederived in Section~\ref{sec:metagame} for completeness.

\begin{proof}[Proof of Theorem~\ref{thm:intro-main}]
Fix $\varepsilon>0$. By Theorem~\ref{thm:meta-equilibrium}, choose a stationary mixed profile
\[
x^*=(x_1^*,\dots,x_n^*)\in X
\]
such that each $x_i^*$ is a best response to $x_{-i}^*$ in the average-reward meta-game.

Now fix a player $i$. Consider the average-reward MDP faced by player $i$ against the stationary mixed opponents' profile $x_{-i}^*$. Let
\[
g_i^*:=g_i^*(x_{-i}^*),\qquad h_i^*:=h_i^{x_{-i}^*},\qquad H_i:=\spn(h_i^*).
\]
Because $x_i^*\in B_i(x_{-i}^*)$, we have
\[
g_i(x^*)=\max_{\mu_i\in X_i}g_i(\mu_i,x_{-i}^*)=g_i^*.
\]
Moreover, by Proposition~\ref{prop:best-response-structure}, the support of $x_i^*(\cdot\mid s)$ is contained in the maximizing set $A_i^*(s;x_{-i}^*)$ for every state $s$.

Let $s\in S$ be any initial state, let $T\ge1$, and let $\tau_i$ be any behavioral deviation of player $i$.

\smallskip
\noindent
\textbf{Step 1: upper bound on the deviator.}
By Proposition~\ref{prop:supermartingale},
\begin{equation}\label{eq:proof-dev-upper}
\gamma_i^T\bigl(s,(\tau_i,x_{-i}^*)\bigr)\le g_i^*+\frac{H_i}{T}.
\end{equation}
This bound holds for \emph{every} behavioral deviation $\tau_i$, including arbitrary history-dependent randomized deviations.

\smallskip
\noindent
\textbf{Step 2: lower bound on the compliant payoff.}
Since $x_i^*$ mixes only over maximizing actions, Proposition~\ref{prop:optimal-stationary}(a) yields the exact identity
\[
\gamma_i^T(s,x^*)=g_i^*+\frac{h_i^*(s)-\E_s^{x^*}[h_i^*(s_{T+1})]}{T}.
\]
Hence
\begin{equation}\label{eq:proof-compliant-lower}
\gamma_i^T(s,x^*)\ge g_i^* - \frac{H_i}{T}.
\end{equation}
By Proposition~\ref{prop:ergodic}, one also has almost sure and $L^1$ convergence
\[
\frac1T\sum_{t=1}^{T}u_i(s_t,a_t)\longrightarrow g_i(x^*)=g_i^*
\qquad\text{under }\Prob_s^{x^*}.
\]
Thus the compliant path has the correct ergodic average, and \eqref{eq:proof-compliant-lower} gives the quantitative finite-horizon estimate needed for uniform equilibrium.

\smallskip
\noindent
\textbf{Step 3: regret bound.}
Subtracting \eqref{eq:proof-compliant-lower} from \eqref{eq:proof-dev-upper} gives
\[
\gamma_i^T\bigl(s,(\tau_i,x_{-i}^*)\bigr)-\gamma_i^T(s,x^*)\le \frac{2H_i}{T}.
\]
Since this holds for every player $i$, every initial state $s$, every horizon $T\ge1$, and every behavioral deviation $\tau_i$, we obtain the uniform estimate
\[
\sup_{i\in N}\sup_{s\in S}\sup_{\tau_i}\Bigl[\gamma_i^T\bigl(s,(\tau_i,x_{-i}^*)\bigr)-\gamma_i^T(s,x^*)\Bigr]
\le \frac{2H}{T},
\]
where
\[
H:=\max_{i\in N}H_i=\max_{i\in N}\spn(h_i^*).
\]

\smallskip
\noindent
\textbf{Step 4: choose the threshold.}
Set
\[
T_0:=\left\lceil \frac{2H}{\varepsilon}\right\rceil.
\]
Then for every $T\ge T_0$,
\[
\gamma_i^T(s,x^*)\ge \gamma_i^T\bigl(s,(\tau_i,x_{-i}^*)\bigr)-\varepsilon
\]
for every player $i$, every initial state $s$, and every behavioral deviation $\tau_i$.

Therefore $x^*$ is a uniform $\varepsilon$-equilibrium. This proves Theorem~\ref{thm:intro-main}.
\end{proof}

For later reference, we record the quantitative consequence explicitly.

\begin{corollary}[Uniform $O(1/T)$ regret]\label{cor:o1overT}
Under Condition~\ref{cond:A}, the stationary profile $x^*$ from Theorem~\ref{thm:meta-equilibrium} satisfies
\[
\gamma_i^T\bigl(s,(\tau_i,x_{-i}^*)\bigr)-\gamma_i^T(s,x^*)
\le \frac{2\spn(h_i^*)}{T}
\le \frac{2\max_j\spn(h_j^*)}{T}
\]
for every player $i$, every initial state $s$, every horizon $T\ge1$, and every behavioral deviation $\tau_i$.
\end{corollary}




\part{Irreducibility despite every player}

\section{The Fudenberg-Yamamoto condition and the guardian profile}\label{sec:FY}

Part~II replaces classical irreducibility by the weaker condition introduced by Fudenberg and Yamamoto \cite{FudenbergYamamoto2011}. In the present perfect-monitoring setup there is no separate public-signal variable, so their definition simplifies as follows.

\begin{condition}[Fudenberg-Yamamoto irreducibility despite player $i$]\label{cond:FY}
A finite stochastic game is \emph{irreducible despite player $i$} if for every ordered pair of states $(s,q)\in S\times S$ there exist an integer $m\ge1$, a sequence of states
\[
s=s_1,s_2,\dots,s_m=q,
\]
and, for each $t<m$, a mixed action profile $\alpha_{-i}^t\in \Delta(A_{-i}(s_t))$ of the opponents such that
\[
\sum_{a_{-i}\in A_{-i}(s_t)}\alpha_{-i}^t(a_{-i})\,P(s_{t+1}\mid s_t,a_i,a_{-i})>0
\qquad\text{for every }a_i\in A_i(s_t).
\]
The game satisfies \emph{Condition~(FY-a)} if it is irreducible despite every player $i$.
\end{condition}

The next proposition is the working form of Condition~\ref{cond:FY}. It is exactly the bounded-return statement extracted from Fudenberg and Yamamoto \cite[Lemma~1]{FudenbergYamamoto2011}. We keep the proof short because the only nontrivial input is their lemma.

For each player $j$ and state $s\in S$, let $u_j(\cdot\mid s)$ denote the uniform distribution on $A_j(s)$, and write
\[
u=(u_j)_{j\in N},\qquad u_{-i}=(u_j)_{j\neq i}.
\]

\begin{proposition}[Uniform guardian profile and bounded returns]\label{prop:guardian}
Assume Condition~\ref{cond:FY}. Fix a player $i$. Then the opponents' statewise uniform profile $u_{-i}$ has the following property. There exists a constant $p_i\in(0,1]$ such that for every initial state $s\in S$, every target state $q\in S$, and every behavioral strategy $\tau_i$ of player $i$,
\begin{equation}\label{eq:guardian-prob}
\Prob_s^{(\tau_i,u_{-i})}\!\bigl(\tau_q\le |S|\bigr)\ge p_i,
\end{equation}
where $\tau_q:=\inf\{t\ge1:s_t=q\}$.

Consequently:
\begin{enumerate}[label=\textup{(\alph*)},leftmargin=2.5em]
\item for every pure stationary strategy $f_i$ of player $i$, the chain $P^{f_i,u_{-i}}$ on $S$ is irreducible;
\item for every target state $q\in S$,
\begin{equation}\label{eq:guardian-mean}
\sup_{s\in S}\sup_{\tau_i}\E_s^{(\tau_i,u_{-i})}[\tau_q]\le \frac{|S|}{p_i}.
\end{equation}
\end{enumerate}
\end{proposition}

\begin{proof}
Fudenberg and Yamamoto \cite[Lemma~1]{FudenbergYamamoto2011} prove that under Condition~\ref{cond:FY}, for each player $i$, the opponents may be chosen to use statewise uniform randomization and there is a constant $p_i>0$ such that, for every ordered pair of states $(s,q)$, state $q$ is reached from $s$ within at most $|S|$ periods with probability at least $p_i$ under every Markov strategy of player $i$. Their proof actually produces a path whose one-step transitions are positive for \emph{every} action of player $i$ at each visited state, so the same lower bound extends immediately to arbitrary behavioral controls by conditioning on the realized action of player $i$ at each stage. This gives \eqref{eq:guardian-prob}.

Statement \textup{(a)} follows immediately, because \eqref{eq:guardian-prob} implies that from every initial state $s$ there is positive probability of reaching every target $q$ in at most $|S|$ steps under $(f_i,u_{-i})$.

For \textup{(b)}, partition time into blocks of length $|S|$. Fix a target $q$. Conditional on the public history at the beginning of any block, the current state is some random state $z\in S$, and by \eqref{eq:guardian-prob} the conditional probability of hitting $q$ within that block is at least $p_i$. Therefore $\tau_q$ is stochastically dominated by $|S|$ times a geometric random variable with success probability $p_i$, which yields \eqref{eq:guardian-mean}.
\end{proof}

The quantifier pattern is genuinely weaker than classical irreducibility.

\begin{example}[Strict weakness of Condition~\ref{cond:FY}]\label{ex:strictly-weaker}
Consider a two-player game with state space $S=\{0,1\}$ and action sets
\[
A_1(0)=A_1(1)=A_2(0)=A_2(1)=\{0,1\}.
\]
Let the transition rule be
\[
s_{t+1}=s_t\oplus a_{1,t}\oplus a_{2,t},
\]
where $\oplus$ is addition modulo $2$, and let stage payoffs be arbitrary bounded numbers, say identically zero. If player~$2$ uses the uniform mixed action $u_2$ at both states, then regardless of player~$1$'s action choice, the next state is $0$ or $1$ with probability $1/2$ each. Thus the chain under $(f_1,u_2)$ is irreducible for every pure stationary $f_1$. By symmetry the same holds for player~$2$, so Condition~\ref{cond:FY} is satisfied.

However classical irreducibility fails. If both players choose action $0$ at both states, then
\[
s_{t+1}=s_t,
\]
so the chain has two absorbing states and is reducible. Hence Condition~\ref{cond:FY} is strictly weaker than classical irreducibility.
\end{example}

For the rest of Part~II we fix:
\begin{itemize}[leftmargin=2.4em]
\item an anchor state $s^\bullet\in S$,
\item the corresponding gate state
\[
\omega^\bullet:=(s^\bullet,1),
\]
\item the constants $p_i$ from Proposition~\ref{prop:guardian}, and
\item the return-phase constants
\[
M_i:=\frac{|S|}{p_i},
\qquad
M:=\max_{i\in N}M_i.
\]
\end{itemize}

\section{Return-to-anchor augmentation}\label{sec:augmentation}

Fix an integer $K\ge1$. Let
\[
M_K:=\{\mathsf R,1,\dots,K\}
\]
be the mode set, where $\mathsf R$ is the return mode and $1,\dots,K$ are work modes. The augmented state space is
\[
\Omega_K:=S\times M_K.
\]
Write a generic augmented state as $\omega=(s,m)$ with physical state $s\in S$ and mode $m\in M_K$.

The mode update is deterministic. Given current augmented state $(s,m)$, current action profile $a\in A(s)$, and next physical state $s'$, define the next mode by
\[
\Phi_K(m,s')=
\begin{cases}
m+1, & \text{if }m\in\{1,\dots,K-1\},\\[0.3em]
\mathsf R, & \text{if }m=K,\\[0.3em]
1, & \text{if }m=\mathsf R\text{ and }s'=s^\bullet,\\[0.3em]
\mathsf R, & \text{if }m=\mathsf R\text{ and }s'\neq s^\bullet.
\end{cases}
\]
Accordingly, the augmented transition kernel is
\[
\widehat P_K\bigl((s',m')\mid (s,m),a\bigr)
=
P(s'\mid s,a)\,\1_{\{m'=\Phi_K(m,s')\}}.
\]
The augmented stage payoff is simply
\[
\widehat u_i\bigl((s,m),a\bigr):=u_i(s,a).
\]

Thus the return mode is regenerative: whenever the process is in mode $\mathsf R$, all prescribed play will be uniform randomization, and the mode remains $\mathsf R$ until the physical state hits $s^\bullet$. The next stage then begins in the gate state $\omega^\bullet=(s^\bullet,1)$. After that there are exactly $K$ work stages, after which the mode returns to $\mathsf R$.

Let
\[
Y_i^K:=\prod_{(s,m)\in\Omega_K}\Delta(A_i(s))
\]
be the full stationary mixed strategy space of player $i$ on the augmented game. The \emph{restricted} stationary strategy space is
\[
X_i^K:=\Bigl\{x_i\in Y_i^K:\ x_i(\cdot\mid (s,\mathsf R))=u_i(\cdot\mid s)\text{ for every }s\in S\Bigr\},
\]
and
\[
X^K:=\prod_{i\in N}X_i^K.
\]
A profile $x^K\in X^K$ is therefore free only in the work modes.

Every $x^K\in X^K$ induces a public finite-memory strategy in the original stochastic game. The public memory variable is the mode $m_t\in M_K$. Start with
\[
m_1:=\mathsf R.
\]
At stage $t$, after observing the current physical state $s_t$ and mode $m_t$, player $i$ chooses action according to
\[
x_i^K(\cdot\mid (s_t,m_t)).
\]
After the next physical state $s_{t+1}$ is observed, the memory updates to
\[
m_{t+1}:=\Phi_K(m_t,s_{t+1}).
\]
Because the mode update depends only on the publicly observed physical-state path, the mode is public. Thus every behavioral deviation in the original game is a behavioral control in the augmented unilateral MDP, and the average payoff in the original game coincides with the average payoff in the augmented game.

The whole point of the augmentation is the following lemma.

\begin{lemma}[Bounded hitting time to the gate state]\label{lem:gate-hit}
Fix a player $i$, fix $x_{-i}\in X_{-i}^K$, and consider either of the following unilateral control problems on $\Omega_K$:
\begin{enumerate}[label=\textup{(\roman*)},leftmargin=2.5em]
\item the \emph{unrestricted} MDP, in which player $i$ may choose any action in every augmented state;
\item the \emph{restricted} MDP, in which player $i$ is also forced to use $u_i(\cdot\mid s)$ in return states $(s,\mathsf R)$ and is free only in work states.
\end{enumerate}
Then for every initial augmented state $\omega\in\Omega_K$ and every behavioral control of player $i$,
\begin{equation}\label{eq:gate-prob}
\Prob_{\omega}\!\bigl(\tau_{\omega^\bullet}\le K+|S|\bigr)\ge p_i,
\end{equation}
where $\tau_{\omega^\bullet}:=\inf\{t\ge1:\omega_t=\omega^\bullet\}$.
Consequently,
\begin{equation}\label{eq:gate-mean}
\sup_{\omega\in\Omega_K}\sup_{\tau_i}\E_{\omega}[\tau_{\omega^\bullet}]
\le \frac{K+|S|}{p_i}.
\end{equation}
\end{lemma}

\begin{proof}
If the initial mode is $\mathsf R$, the process stays in return mode until the physical state hits $s^\bullet$, and the opponents of player $i$ use the guardian profile $u_{-i}$ during that entire episode. Proposition~\ref{prop:guardian} therefore gives
\[
\Prob_{\omega}\!\bigl(\tau_{\omega^\bullet}\le |S|\bigr)\ge p_i.
\]

If the initial mode is $m\in\{1,\dots,K\}$, then after at most $K-m+1\le K$ stages the process enters a return state. Conditional on the physical state at that time, Proposition~\ref{prop:guardian} again gives probability at least $p_i$ of hitting $s^\bullet$, hence $\omega^\bullet$, within the next $|S|$ stages. This proves \eqref{eq:gate-prob}.

For \eqref{eq:gate-mean}, partition time into blocks of length $K+|S|$. Conditional on the public history at the beginning of any block, the current augmented state is some random $\omega$, and by \eqref{eq:gate-prob} the conditional probability of hitting $\omega^\bullet$ within that block is at least $p_i$. Therefore $\tau_{\omega^\bullet}$ is stochastically dominated by $(K+|S|)$ times a geometric random variable with success probability $p_i$.
\end{proof}

The bounded-return lemma collapses the recurrent-class problem on the augmented space.

\begin{proposition}[Unichain lemma on the augmented space]\label{prop:aug-unichain}
Fix a player $i$ and $x_{-i}\in X_{-i}^K$.
\begin{enumerate}[label=\textup{(\alph*)},leftmargin=2.5em]
\item Every pure stationary policy of the unrestricted unilateral MDP induces a Markov chain on $\Omega_K$ with exactly one recurrent class.
\item Every pure stationary policy of the restricted unilateral MDP induces a Markov chain on $\Omega_K$ with exactly one recurrent class.
\end{enumerate}
In both cases the unique recurrent class contains the gate state $\omega^\bullet=(s^\bullet,1)$.
\end{proposition}

\begin{proof}
Fix any pure stationary policy in either MDP. By Lemma~\ref{lem:gate-hit}, starting from every augmented state there is positive probability of reaching $\omega^\bullet$ within $K+|S|$ steps. Hence every recurrent class must contain $\omega^\bullet$. Since a finite Markov chain has at least one recurrent class, there can be only one.
\end{proof}

\begin{corollary}\label{cor:aug-profile-unichain}
For every restricted stationary profile $x\in X^K$, the induced Markov chain on $\Omega_K$ has exactly one recurrent class and therefore a unique invariant distribution.
\end{corollary}

\begin{proof}
Fix any player $i$. Under the profile $x$, the behavior of player $i$ is a particular behavioral control in the unrestricted augmented MDP against $x_{-i}$. By Lemma~\ref{lem:gate-hit}, from every augmented state there is positive probability of reaching $\omega^\bullet$ within $K+|S|$ stages. Therefore every recurrent class of $\widehat P_K^x$ contains $\omega^\bullet$, so the recurrent class is unique.
\end{proof}

At this point the only MDP input needed from Part~I is the following observation.

\begin{remark}[Why the Part~I machinery applies on $\Omega_K$]\label{rem:verbatim}
The proofs of Theorem~\ref{thm:aroe}, Proposition~\ref{prop:supermartingale}, and Proposition~\ref{prop:optimal-stationary} use only finite-state unichain MDP theory. They do not use irreducibility of the \emph{full} state space beyond the unichain consequence. Hence, once Proposition~\ref{prop:aug-unichain} is proved, those statements apply verbatim to the unrestricted and restricted unilateral MDPs on $\Omega_K$.
\end{remark}

\section{The restricted augmented meta-game}\label{sec:aug-metagame}

Fix $K\ge1$. For a restricted stationary profile $x\in X^K$, define the expected one-step payoff of player $i$ in augmented state $\omega=(s,m)$ by
\[
\overline u_i^{K,x}(\omega):=\sum_{a\in A(s)}x(a\mid \omega)\,\widehat u_i(\omega,a)
=
\sum_{a\in A(s)}x(a\mid \omega)\,u_i(s,a).
\]
Let $\widehat P_K^x$ be the induced transition matrix on $\Omega_K$, and let $\widehat\pi^{K,x}$ be its unique invariant distribution, whose existence follows from Corollary~\ref{cor:aug-profile-unichain}. Define the augmented average-reward payoff of player $i$ by
\begin{equation}\label{eq:aug-payoff}
g_i^K(x):=\sum_{\omega\in\Omega_K}\widehat\pi^{K,x}(\omega)\,\overline u_i^{K,x}(\omega).
\end{equation}

\begin{proposition}[Continuity on the restricted augmented space]\label{prop:aug-continuity}
For every player $i$, the maps
\[
x\longmapsto \widehat\pi^{K,x}\in\Delta(\Omega_K)
\qquad\text{and}\qquad
x\longmapsto g_i^K(x)
\]
are continuous on $X^K$.
\end{proposition}

\begin{proof}
The map $x\mapsto \widehat P_K^x$ is multilinear, hence continuous. Let $x^m\to x$ in $X^K$ and write $\widehat\pi^m:=\widehat\pi^{K,x^m}$. Every convergent subsequence of $(\widehat\pi^m)$ has a limit point $\bar\pi\in\Delta(\Omega_K)$ satisfying
\[
\bar\pi\,\widehat P_K^x=\bar\pi.
\]
By Corollary~\ref{cor:aug-profile-unichain}, $\widehat P_K^x$ has a unique invariant distribution, so $\bar\pi=\widehat\pi^{K,x}$. Hence the whole sequence converges. Continuity of $g_i^K$ follows from \eqref{eq:aug-payoff}.
\end{proof}

Fix a player $i$ and an opponents' restricted stationary profile $x_{-i}\in X_{-i}^K$. Let
\[
g_i^{K,R}(x_{-i})
\]
be the optimal average reward in the restricted unilateral MDP on $\Omega_K$, and let
\[
g_i^{K,U}(x_{-i})
\]
be the optimal average reward in the unrestricted unilateral MDP. By Proposition~\ref{prop:aug-unichain} and Remark~\ref{rem:verbatim}, both values are attained by stationary policies and admit normalized biases at the reference state $\omega^\bullet$.

The restricted best-response correspondence is
\[
B_i^K(x_{-i}):=\arg\max_{y_i\in X_i^K}g_i^K(y_i,x_{-i}).
\]

\begin{proposition}[Structure of restricted best responses]\label{prop:aug-br-structure}
Fix $x_{-i}\in X_{-i}^K$. Let $A_{i,K}^{*}(\omega;x_{-i})$ be the maximizing action set in the restricted average-reward optimality equation for player $i$ at the work state $\omega\in S\times\{1,\dots,K\}$. Then
\[
B_i^K(x_{-i})
=
\Bigl\{y_i\in X_i^K:\ \supp y_i(\cdot\mid \omega)\subseteq A_{i,K}^{*}(\omega;x_{-i})
\text{ for every }\omega\in S\times\{1,\dots,K\}\Bigr\}.
\]
In particular, $B_i^K(x_{-i})$ is nonempty, compact, and convex.
\end{proposition}

\begin{proof}
This is the same argument as Proposition~\ref{prop:best-response-structure}. On return states there is no choice because $y_i$ must equal the uniform distribution. On work states, Remark~\ref{rem:verbatim} reduces the problem to the finite unichain MDP characterization of optimal stationary responses.
\end{proof}

\begin{proposition}[Closed graph of the restricted best-response correspondence]\label{prop:aug-br-closed}
For every player $i$, the correspondence $B_i^K:X_{-i}^K\rightrightarrows X_i^K$ has closed graph and therefore is upper hemicontinuous.
\end{proposition}

\begin{proof}
The proof is identical to Proposition~\ref{prop:best-response-closed-graph}, using Proposition~\ref{prop:aug-continuity} in place of Proposition~\ref{prop:continuity}.
\end{proof}

\begin{theorem}[Restricted stationary equilibrium on $\Omega_K$]\label{thm:aug-eq}
For every $K\ge1$ there exists a restricted stationary profile
\[
x^{K,*}\in X^K
\]
such that for every player $i$,
\[
g_i^K(x^{K,*})\ge g_i^K(y_i,x_{-i}^{K,*})\qquad\text{for every }y_i\in X_i^K.
\]
Equivalently,
\[
x_i^{K,*}\in B_i^K(x_{-i}^{K,*})\qquad\text{for every }i.
\]
\end{theorem}

\begin{proof}
Define $B^K:X^K\rightrightarrows X^K$ by
\[
B^K(x):=\prod_{i\in N}B_i^K(x_{-i}).
\]
The domain $X^K$ is a nonempty compact convex subset of a Euclidean space. By Propositions~\ref{prop:aug-br-structure} and \ref{prop:aug-br-closed}, the correspondence $B^K$ has nonempty compact convex values and closed graph. Therefore Theorem~\ref{thm:kfg} yields a fixed point $x^{K,*}\in B^K(x^{K,*})$.
\end{proof}

\section{Regenerative comparison and bias bounds}\label{sec:regenerative}

This section contains the one genuinely new estimate relative to Part~I: the unrestricted and restricted average rewards differ only on the return phase, so their gains differ by $O(1/K)$.

Fix a player $i$, fix $K\ge1$, and fix the opponents' restricted stationary profile $x_{-i}\in X_{-i}^K$. For a stationary policy pair $(\alpha,\beta)$ of player $i$, write
\[
g_i^K(\alpha,\beta;x_{-i})
\]
for the resulting average reward of player $i$ in the unrestricted augmented MDP against $x_{-i}$.

A stationary policy of player $i$ in the unrestricted augmented MDP can be decomposed as a pair $(\alpha,\beta)$, where
\begin{itemize}[leftmargin=2.4em]
\item $\alpha$ is the work-mode component, namely a stationary mixed action at each $\omega\in S\times\{1,\dots,K\}$,
\item $\beta$ is the return-mode component, namely a stationary mixed action at each $\omega\in S\times\{\mathsf R\}$.
\end{itemize}
The restricted policies are exactly the pairs $(\alpha,u_i)$.

Start the augmented process at the gate state $\omega^\bullet=(s^\bullet,1)$. Because of Proposition~\ref{prop:aug-unichain}, $\omega^\bullet$ is recurrent under every stationary policy. Let
\[
\theta_0:=0,
\qquad
\theta_{r+1}:=\inf\{t>\theta_r:\omega_t=\omega^\bullet\}
\]
be the successive return times to $\omega^\bullet$.

\begin{lemma}[Cycle representation]\label{lem:cycle}
For every stationary policy $(\alpha,\beta)$ of player $i$:
\begin{enumerate}[label=\textup{(\alph*)},leftmargin=2.5em]
\item the cycles
\[
\bigl((\omega_t,a_t)\bigr)_{\theta_r<t\le \theta_{r+1}},\qquad r\ge0,
\]
are i.i.d. under $\Prob_{\omega^\bullet}^{(\alpha,\beta),x_{-i}}$;
\item each cycle consists of exactly $K$ work stages followed by a return phase;
\item if
\[
W_i^K(\alpha):=\E_{\omega^\bullet}^{(\alpha,\beta),x_{-i}}
\Bigl[\sum_{t=1}^{K}\widehat u_i(\omega_t,a_t)\Bigr],
\]
\[
B_i^K(\alpha,\beta):=
\E_{\omega^\bullet}^{(\alpha,\beta),x_{-i}}
\Bigl[\sum_{t=K+1}^{\theta_1}\widehat u_i(\omega_t,a_t)\Bigr],
\]
and
\[
L_i^K(\alpha,\beta):=
\E_{\omega^\bullet}^{(\alpha,\beta),x_{-i}}[\theta_1-K],
\]
then
\begin{equation}\label{eq:cycle-ratio}
g_i^K(\alpha,\beta;x_{-i})
=
\frac{W_i^K(\alpha)+B_i^K(\alpha,\beta)}{K+L_i^K(\alpha,\beta)}.
\end{equation}
Moreover,
\begin{equation}\label{eq:return-bounds}
0\le W_i^K(\alpha)\le K,
\qquad
0\le B_i^K(\alpha,\beta)\le L_i^K(\alpha,\beta)\le M_i.
\end{equation}
\end{enumerate}
\end{lemma}

\begin{proof}
The strong Markov property at the recurrent state $\omega^\bullet$ gives \textup{(a)}. Statement \textup{(b)} is built into the mode dynamics: after entering $\omega^\bullet$ the process deterministically traverses work modes $1,\dots,K$, and only then can it enter return mode. Formula \eqref{eq:cycle-ratio} is the standard renewal-reward ratio for regenerative processes with finite mean cycle length.

For \eqref{eq:return-bounds}, the bound on $W_i^K(\alpha)$ is immediate from the payoff normalization $u_i\in[0,1]$. Also $B_i^K(\alpha,\beta)\le L_i^K(\alpha,\beta)$ because each return-mode stage payoff is at most $1$. Finally, once the process enters return mode, Proposition~\ref{prop:guardian} implies that the expected number of return-mode stages before hitting $s^\bullet$, hence before re-entering $\omega^\bullet$, is at most $M_i=|S|/p_i$.
\end{proof}

The ratio \eqref{eq:cycle-ratio} is close to the work-block average $W_i^K(\alpha)/K$ uniformly in the return policy.

\begin{lemma}[Comparison with the work-block average]\label{lem:work-baseline}
For every stationary policy pair $(\alpha,\beta)$,
\begin{equation}\label{eq:baseline-bound}
\left|
g_i^K(\alpha,\beta;x_{-i})-\frac{W_i^K(\alpha)}{K}
\right|
\le \frac{M_i}{K}.
\end{equation}
\end{lemma}

\begin{proof}
Using \eqref{eq:cycle-ratio},
\[
g_i^K(\alpha,\beta;x_{-i})-\frac{W_i^K(\alpha)}{K}
=
\frac{K\,B_i^K(\alpha,\beta)-W_i^K(\alpha)\,L_i^K(\alpha,\beta)}
{K\bigl(K+L_i^K(\alpha,\beta)\bigr)}.
\]
By \eqref{eq:return-bounds},
\[
0\le B_i^K(\alpha,\beta)\le L_i^K(\alpha,\beta)\le M_i
\qquad\text{and}\qquad
0\le W_i^K(\alpha)\le K.
\]
Hence the numerator has absolute value at most $K\,L_i^K(\alpha,\beta)\le KM_i$, proving \eqref{eq:baseline-bound}.
\end{proof}

\begin{proposition}[The $O(1/K)$ gain comparison]\label{prop:gain-comparison}
For every player $i$, every $K\ge1$, and every opponents' restricted stationary profile $x_{-i}\in X_{-i}^K$,
\[
0\le g_i^{K,U}(x_{-i})-g_i^{K,R}(x_{-i})\le \frac{2M_i}{K}.
\]
\end{proposition}

\begin{proof}
Let
\[
w_i^K(x_{-i}):=\sup_{\alpha}\frac{W_i^K(\alpha)}{K},
\]
where the supremum runs over stationary work-mode policies $\alpha$ of player $i$.

By Lemma~\ref{lem:work-baseline},
\[
g_i^K(\alpha,\beta;x_{-i})\le \frac{W_i^K(\alpha)}{K}+\frac{M_i}{K}\le w_i^K(x_{-i})+\frac{M_i}{K}
\]
for every unrestricted stationary pair $(\alpha,\beta)$. Taking the supremum over $(\alpha,\beta)$ gives
\[
g_i^{K,U}(x_{-i})\le w_i^K(x_{-i})+\frac{M_i}{K}.
\]

Similarly, for every work-mode policy $\alpha$,
\[
g_i^K(\alpha,u_i;x_{-i})\ge \frac{W_i^K(\alpha)}{K}-\frac{M_i}{K}.
\]
Taking the supremum over $\alpha$ yields
\[
g_i^{K,R}(x_{-i})\ge w_i^K(x_{-i})-\frac{M_i}{K}.
\]
Subtracting the two inequalities gives the bound. The lower bound $g_i^{K,U}\ge g_i^{K,R}$ is immediate because the unrestricted action set contains the restricted one.
\end{proof}

We also need a growth bound for the biases on the augmented space.

\begin{proposition}[Bias-span bound on the augmented space]\label{prop:aug-bias-span}
Fix a player $i$, $K\ge1$, and $x_{-i}\in X_{-i}^K$. Let $h$ be a normalized bias at the reference state $\omega^\bullet$ for either the unrestricted or the restricted average-reward optimality equation on $\Omega_K$. Then
\[
\spn(h)\le 2D_{i,K},
\qquad
D_{i,K}:=\frac{K+|S|}{p_i}.
\]
Consequently, if
\[
H_i(K):=\spn(h_i^{K,U})+\spn(h_i^{K,R}),
\]
then
\[
H_i(K)\le 4D_{i,K}=4\,\frac{K+|S|}{p_i}.
\]
In particular, with
\[
H(K):=\max_{i\in N}H_i(K),
\]
one has $H(K)=O(K)$.
\end{proposition}

\begin{proof}
Fix one of the two MDPs and let $f$ be an optimal stationary policy that selects Bellman maximizers. By Remark~\ref{rem:verbatim} and Proposition~\ref{prop:optimal-stationary}, the process
\[
M_t:=h(\omega_t)+\sum_{u=1}^{t-1}\bigl(\widehat u_i(\omega_u,a_u)-g\bigr)
\]
is a martingale under $\Prob_{\omega}^{f,x_{-i}}$, where $g$ is the corresponding optimal average reward. Let
\[
\tau^\bullet:=\inf\{t\ge1:\omega_t=\omega^\bullet\}.
\]
By Lemma~\ref{lem:gate-hit}, $\E_{\omega}^{f,x_{-i}}[\tau^\bullet]\le D_{i,K}$ for every $\omega$.

Optional stopping at $\tau^\bullet\wedge n$ and then letting $n\to\infty$ gives
\[
h(\omega)
=
\E_{\omega}^{f,x_{-i}}
\Bigl[\sum_{t=1}^{\tau^\bullet-1}\bigl(\widehat u_i(\omega_t,a_t)-g\bigr)\Bigr],
\]
because $h(\omega^\bullet)=0$ by normalization and because $\tau^\bullet$ has finite mean. Since both $\widehat u_i$ and $g$ lie in $[0,1]$, each summand is in $[-1,1]$. Therefore
\[
|h(\omega)|\le \E_{\omega}^{f,x_{-i}}[\tau^\bullet-1]\le D_{i,K}.
\]
Taking the difference between two states yields $\spn(h)\le 2D_{i,K}$. The bound on $H_i(K)$ follows by adding the unrestricted and restricted estimates.
\end{proof}

\section{Proof of the Fudenberg-Yamamoto theorem}\label{sec:proof-fy}

We now prove Theorem~\ref{thm:intro-fy}. Fix $K\ge1$ and choose the restricted stationary equilibrium
\[
x^{K,*}\in X^K
\]
from Theorem~\ref{thm:aug-eq}. Let $\sigma^K$ be the induced public finite-memory strategy profile in the original stochastic game.

Fix a player $i$. Against $x_{-i}^{K,*}$ there are two unilateral MDPs on the augmented state space:
\begin{itemize}[leftmargin=2.4em]
\item the unrestricted MDP, with optimal gain $g_i^{K,U}$ and normalized bias $h_i^{K,U}$;
\item the restricted MDP, with optimal gain $g_i^{K,R}$ and normalized bias $h_i^{K,R}$.
\end{itemize}
Both biases are normalized at $\omega^\bullet$.

\begin{proof}[Proof of Theorem~\ref{thm:intro-fy}]
Fix an initial physical state $s\in S$ and let the augmented process start from
\[
\omega_1:=(s,\mathsf R).
\]
Let $\widehat\gamma_i^T$ denote the $T$-stage average payoff in the augmented game. Because the augmented stage payoffs are exactly the original stage payoffs and the mode is a public function of the physical history, one has
\[
\widehat\gamma_i^T(\omega_1,\cdot)=\gamma_i^T(s,\cdot)
\]
for the induced profile $\sigma^K$ and for every behavioral deviation in the original game.

Now fix an arbitrary behavioral deviation $\tau_i$ of player $i$ in the original game, viewed equivalently as a behavioral control in the unrestricted augmented MDP.

\smallskip
\noindent
\textbf{Step 1: upper bound on the deviator.}
By Remark~\ref{rem:verbatim} and Proposition~\ref{prop:supermartingale}, applied to the unrestricted augmented MDP,
\begin{equation}\label{eq:fy-dev-upper}
\widehat\gamma_i^T\bigl(\omega_1,(\tau_i,x_{-i}^{K,*})\bigr)
\le
g_i^{K,U}+\frac{\spn(h_i^{K,U})}{T}.
\end{equation}

\smallskip
\noindent
\textbf{Step 2: lower bound on the compliant payoff.}
Because $x_i^{K,*}$ is a restricted stationary best response to $x_{-i}^{K,*}$, the profile $(x_i^{K,*},x_{-i}^{K,*})$ attains the restricted optimal gain $g_i^{K,R}$. By Remark~\ref{rem:verbatim} and Proposition~\ref{prop:optimal-stationary}, applied to the restricted augmented MDP,
\begin{equation}\label{eq:fy-compliant-lower}
\widehat\gamma_i^T(\omega_1,x^{K,*})
\ge
g_i^{K,R}-\frac{\spn(h_i^{K,R})}{T}.
\end{equation}

\smallskip
\noindent
\textbf{Step 3: compare the two gains.}
By Proposition~\ref{prop:gain-comparison},
\[
g_i^{K,U}-g_i^{K,R}\le \frac{2M_i}{K}.
\]
Subtracting \eqref{eq:fy-compliant-lower} from \eqref{eq:fy-dev-upper} therefore yields
\[
\widehat\gamma_i^T\bigl(\omega_1,(\tau_i,x_{-i}^{K,*})\bigr)-\widehat\gamma_i^T(\omega_1,x^{K,*})
\le
\frac{2M_i}{K}+\frac{H_i(K)}{T},
\]
where
\[
H_i(K):=\spn(h_i^{K,U})+\spn(h_i^{K,R}).
\]

\smallskip
\noindent
\textbf{Step 4: return to the original game.}
Because augmented and original average payoffs coincide under the induced finite-memory implementation,
\[
\gamma_i^T\bigl(s,(\tau_i,\sigma_{-i}^K)\bigr)-\gamma_i^T(s,\sigma^K)
\le
\frac{2M_i}{K}+\frac{H_i(K)}{T}.
\]
Taking the supremum over $i$, $s$, and $\tau_i$ gives
\[
\sup_{i,s,\tau_i}
\Bigl[\gamma_i^T\bigl(s,(\tau_i,\sigma_{-i}^K)\bigr)-\gamma_i^T(s,\sigma^K)\Bigr]
\le
\frac{2M}{K}+\frac{H(K)}{T},
\]
where
\[
M:=\max_i M_i,
\qquad
H(K):=\max_i H_i(K).
\]

\smallskip
\noindent
\textbf{Step 5: choose the parameters.}
Given $\varepsilon>0$, let
\[
K_\varepsilon:=\left\lceil \frac{4M}{\varepsilon}\right\rceil
\qquad\text{and}\qquad
T_0(\varepsilon):=\left\lceil \frac{2H(K_\varepsilon)}{\varepsilon}\right\rceil.
\]
Then for every $T\ge T_0(\varepsilon)$,
\[
\frac{2M}{K_\varepsilon}+\frac{H(K_\varepsilon)}{T}
\le \varepsilon.
\]
Hence
\[
\gamma_i^T(s,\sigma^{K_\varepsilon})
\ge
\gamma_i^T\bigl(s,(\tau_i,\sigma_{-i}^{K_\varepsilon})\bigr)-\varepsilon
\]
for every player $i$, every initial state $s$, every horizon $T\ge T_0(\varepsilon)$, and every behavioral deviation $\tau_i$.

Therefore $\sigma^{K_\varepsilon}$ is a uniform $\varepsilon$-equilibrium.
\end{proof}

For later reference we isolate the quantitative estimate.

\begin{corollary}[Uniform finite-horizon regret under Condition~\ref{cond:FY}]\label{cor:fy-regret}
For every $K\ge1$, the induced finite-memory profile $\sigma^K$ satisfies
\[
\gamma_i^T\bigl(s,(\tau_i,\sigma_{-i}^K)\bigr)-\gamma_i^T(s,\sigma^K)
\le
\frac{2M_i}{K}+\frac{H_i(K)}{T}
\le
\frac{2M}{K}+\frac{H(K)}{T}
\]
for every player $i$, every initial state $s$, every horizon $T\ge1$, and every behavioral deviation $\tau_i$.
\end{corollary}

\begin{remark}[How Part~II differs from Part~I]
The proof above uses exactly the same Bellman-supermartingale argument as Part~I, but on the augmented state space $\Omega_K$ rather than on the original state space $S$. The only new input is Proposition~\ref{prop:gain-comparison}, which quantifies the cost of forcing the return mode to use the guardian profile. The transient states of the augmented chain require no separate treatment: once Proposition~\ref{prop:aug-unichain} is established, the bias on $\Omega_K$ already accounts for them.
\end{remark}

\section{Discussion}\label{sec:discussion-new}

Part~I is the classical v28 argument. It shows that once the original state space already yields a unichain unilateral control problem under every stationary profile, a stationary average-reward equilibrium automatically yields an $O(1/T)$ uniform-equilibrium bound.

Part~II is the real correction. The Fudenberg-Yamamoto condition is not the universal quantifier from v28's old Condition~A. It is existential in the opponents' profile, and precisely that feature makes it weaker than classical irreducibility. The return-to-anchor augmentation converts that weaker hypothesis into a unichain statement on a larger state space. After the augmentation, the v28 Bellman argument survives unchanged, and the only new estimate is the regenerative $O(1/K)$ comparison between unrestricted and restricted gains.

What remains open is the genuinely general multiplayer case without any guardian profile of the Fudenberg-Yamamoto type. There the recurrent-class selection problem reappears: a deviator may be able to change the long-run weights on several recurrent classes with different gains. The present paper does not address that difficulty. It proves uniform $\varepsilon$-equilibrium only under classical irreducibility or, more generally, under irreducibility despite every player.


\begin{thebibliography}{99}

\bibitem{Federgruen1978}
A. Federgruen.
\newblock On $N$-person stochastic games with denumerable state space.
\newblock \emph{Advances in Applied Probability}, 10(2):452--471, 1978.

\bibitem{FudenbergYamamoto2011}
D. Fudenberg and Y. Yamamoto.
\newblock The folk theorem for irreducible stochastic games with imperfect public monitoring.
\newblock \emph{Journal of Economic Theory}, 146(4):1664--1683, 2011.

\bibitem{Glicksberg1952}
I.~L. Glicksberg.
\newblock A further generalization of the Kakutani fixed point theorem, with application to Nash equilibrium points.
\newblock \emph{Proceedings of the American Mathematical Society}, 3(1):170--174, 1952.

\bibitem{MertensNeyman1981}
J.-F. Mertens and A. Neyman.
\newblock Stochastic games.
\newblock \emph{International Journal of Game Theory}, 10(2):53--66, 1981.

\bibitem{Neyman2003Minmax}
A. Neyman.
\newblock Stochastic games: existence of the minmax.
\newblock In A. Neyman and S. Sorin, editors, \emph{Stochastic Games and Applications}, volume 570 of \emph{NATO Science Series C: Mathematical and Physical Sciences}, pages 173--193. Springer, Dordrecht, 2003.

\bibitem{Neyman2003MarkovChains}
A. Neyman.
\newblock From Markov chains to stochastic games.
\newblock In A. Neyman and S. Sorin, editors, \emph{Stochastic Games and Applications}, volume 570 of \emph{NATO Science Series C: Mathematical and Physical Sciences}, pages 9--25. Springer, Dordrecht, 2003.

\bibitem{Norris1998}
J.~R. Norris.
\newblock \emph{Markov Chains}.
\newblock Cambridge University Press, Cambridge, 1998.

\bibitem{OpenProblems2026}
Game Theory Society (Israel Chapter).
\newblock Open problems.
\newblock \url{https://sites.google.com/view/thegametheorysocietyilchapter/open-problems}, accessed April 2, 2026.

\bibitem{Puterman1994}
M.~L. Puterman.
\newblock \emph{Markov Decision Processes: Discrete Stochastic Dynamic Programming}.
\newblock Wiley, New York, 1994.

\bibitem{Rogers1969}
P.~D. Rogers.
\newblock \emph{Nonzero-Sum Stochastic Games}.
\newblock PhD thesis, University of California, Berkeley, 1969.
\newblock Operations Research Center Report ORC 69--8.

\bibitem{Shapley1953}
L.~S. Shapley.
\newblock Stochastic games.
\newblock \emph{Proceedings of the National Academy of Sciences of the United States of America}, 39(10):1095--1100, 1953.

\bibitem{Sobel1971}
M.~J. Sobel.
\newblock Noncooperative stochastic games.
\newblock \emph{The Annals of Mathematical Statistics}, 42(6):1930--1935, 1971.

\bibitem{Solan2022}
E. Solan.
\newblock \emph{A Course in Stochastic Game Theory}.
\newblock Cambridge University Press, Cambridge, 2022.

\bibitem{SolanVieille2015}
E. Solan and N. Vieille.
\newblock Stochastic games.
\newblock \emph{Proceedings of the National Academy of Sciences of the United States of America}, 112(45):13743--13746, 2015.

\bibitem{Vieille2000a}
N. Vieille.
\newblock Two-player stochastic games {I}: A reduction.
\newblock \emph{Israel Journal of Mathematics}, 119:55--91, 2000.

\bibitem{Vieille2000b}
N. Vieille.
\newblock Two-player stochastic games {II}: The case of recursive games.
\newblock \emph{Israel Journal of Mathematics}, 119:93--126, 2000.

\end{thebibliography}

\end{document}
